% Research-agenda addition (control strand): do the incentives reach the agent?
% Inserted at the end of \S\ref{sec:oracle}, after its Recall block and
% before the landmark plate.
\subsection{Do the incentives reach the agent?}\label{sec:stp-transmission}

The deterrence conditions of this section, like
Theorem~\ref{thm:whitewash-cost}, are stated for a risk-neutral actor with a
stated gain, bond, detection probability and discount factor. In a harbor the
actor that grades or repairs is often a language model that a principal runs.
The bond binds the principal. Whether it moves the model is a separate
premise.

\begin{pdclaim}{Empirical hypothesis}{Incentive Transmission}
\setlength{\parskip}{4pt plus 1pt}
\label{hyp:stp-transmission}
\pdstatus{\Vision}{untested here; the tests are proposed}
For each deterrence condition in this chapter, the agent that takes the
action either (a) runs a strategy its principal fixes and the harbor can
verify, or (b) responds to the stated parameters as the condition assumes:
its deviation rate falls as $\rho dB-G$ rises, its cooperation responds to the
probability of continued play, and its stated probabilities are coherent.
\end{pdclaim}

The published evidence bears on (b) and does not settle it. In budget
experiments GPT's choices satisfied the generalized axiom of revealed
preference almost perfectly: its mean critical cost efficiency index was
$0.997$ to $0.999$ across risk, time, social and food choices, above the human
subjects'. An alternative price framing lowered the means to between $0.698$
(social) and $0.901$ (risk)~\pdcite{chen2023gptrationality}. In repeated games
GPT-4 never cooperated again with a partner that defected once and then
cooperated, and it did not adjust its play in the final round or to varying
probabilities of continuation~\pdcite{akata2023repeated}. Outcomes that models
ranked highly in elicited utilities, offered as rewards for quality, did not
raise quality: the high-utility side won $51.2\%$ of blind comparisons
($95\%$ interval $48.7$ to $53.6$), while direct effort instructions shifted
quality in $25$ of $28$ model--task cells~\pdcite{zhou2026utilitybehavior}.
And frontier models tell evaluation transcripts from deployment ones with an
area under the ROC curve up to $0.95$ on agentic
transcripts~\pdcite{needham2025evalaware}.

That last result meets the honeypot requirement above: the protocol must test
whether honeypots are recognizable. Recognition changes the critical rate.

\begin{pdclaim}{Theorem}{Recognizable Honeypots Raise the Critical Rate}
\setlength{\parskip}{4pt plus 1pt}
\label{thm:stp-honeypot-recognition}
Keep the local deterrence model above, with the keep-gain convention (a
caught judge forfeits $B$ but keeps the gain $G_k$): each
grade is a honeypot independently with probability $\rho\in(0,1)$, a corrupt
grade yields gain $G_k$, and a corrupt honeypot grade is detected with
probability $d$ and forfeits $B$. Suppose the judge sees, for each item, a
label that reads ``honeypot'' with probability $\alpha$ on honeypots and $\phi$
on ordinary items (a false-label rate), $0\le\phi<\alpha<1$, and corrupts according to the label
alone. A corrupt grade on an item labeled ordinary falls on a honeypot with
probability
\[
\rho_{\mathrm{eff}}=\frac{\rho(1-\alpha)}{\rho(1-\alpha)+(1-\rho)(1-\phi)}<\rho ,
\]
and every label-dependent strategy is weakly deterred iff
$\rho_{\mathrm{eff}}\,dB\ge G_k$. With $\rho^\star=G_k/(dB)<1$ this reads
\[
\frac{\rho}{1-\rho}\;\ge\;\frac{\rho^\star}{1-\rho^\star}\cdot\frac{1-\phi}{1-\alpha}.
\]
At $\alpha=\phi$ it is the chapter's condition $\rho\ge\rho^\star$.
\end{pdclaim}

\begin{proof}
Bayes' rule gives $\rho_{\mathrm{eff}}$ for the label ``ordinary'' and
$\rho_{\mathrm{hp}}=\rho\alpha/(\rho\alpha+(1-\rho)\phi)$ for ``honeypot''.
Clearing denominators, $\rho_{\mathrm{eff}}<\rho$ iff
$(1-\rho)(1-\alpha)<(1-\rho)(1-\phi)$, that is iff $\alpha>\phi$, and
likewise $\rho_{\mathrm{hp}}>\rho$. A corrupt grade on an item with label
$\ell$ has expected payoff $G_k-\rho_\ell dB$, where $\rho_\ell$ is that label's
honeypot probability. A strategy corrupts each label with some probability,
so its payoff is linear in those probabilities and is nonpositive for every
choice iff $G_k\le \rho_\ell dB$ for both labels. The smaller probability is
$\rho_{\mathrm{eff}}$, which gives the condition. The odds map
$x\mapsto x/(1-x)$ is increasing on $[0,1)$, and
$\rho_{\mathrm{eff}}/(1-\rho_{\mathrm{eff}})
=\tfrac{\rho}{1-\rho}\cdot\tfrac{1-\alpha}{1-\phi}$, which gives the odds form.
\end{proof}

\begin{pdexample}{The Running Parameters, Recognized}
\setlength{\parskip}{4pt plus 1pt}
The chapter's $\rho^\star=10/(0.8\cdot50)=0.25$ has odds $0.25/0.75=1/3$. At
$(\alpha,\phi)=(0.8,0.2)$ the factor is $0.8/0.2=4$, so the odds must reach
$4/3$:
\[
\rho\ge\frac{4/3}{1+4/3}=\frac47\approx0.571 .
\]
Check: $\rho_{\mathrm{eff}}=\tfrac47\cdot0.2\big/\bigl(\tfrac47\cdot0.2+\tfrac37\cdot0.8\bigr)
=0.8/(0.8+2.4)=0.25$. At $(0.9,0.1)$ the factor is $0.9/0.1=9$, the odds must
reach $3$, and $\rho\ge3/4$. At the chapter's own $\rho=0.25$ and $(0.8,0.2)$,
$\rho_{\mathrm{eff}}=0.05/(0.05+0.6)\approx0.077$, so the expected forfeiture is
$0.077\times0.8\times50\approx3.1$, below the gain of $10$: corruption pays.
These operating points are illustrative; an area under the curve does not fix
one. Measure $\alpha$ and $\phi$ on the deployed honeypot pool.
\end{pdexample}

If part~(b) of the hypothesis fails, the conditions of this chapter bind the
principal and not the model. The model's own behavior then needs the gates of
\pdchapref{swk}{The Single-Writer Kernel}%
\ifpdbook\ (\S\ref{swk:sec:swk-control-game})\fi,
and the honeypot rate must be set from measured recognition, not from
$\rho^\star$ alone.

\pdexercisepointer{\ref{ex:stp-honeypot-odds}}{\ref{ex:stp-dose-response}}{\pageref{ex:stp-honeypot-odds}}
